\documentclass[../../main.tex]{subfiles}
\begin{document}

{\bf [解答]}

% TODO: 概要図

$k=0,1,\cdots,n-1$に対して$L_k=|PQ_k|^2$とおくと，
\begin{align}
  L_k
  &=\left(\cos\dfrac{2\pi}{n}k-a\right)^2+\left(\sin\dfrac{2\pi}{n}k-b\right)^2  \\
  &=1+a^2+b^2-2\left(a\cos\dfrac{2\pi}{n}k+b\sin\dfrac{2\pi}{n}k\right)
\end{align}
となる．$a_k=\left(a\cos\dfrac{2\pi}{n}k+b\sin\dfrac{2\pi}{n}k\right)$とおけば区分求積により，
\begin{align}
  \lim_{n\to\infty}\dfrac{1}{n}\sum_{k=0}^{n-1}a_k
  &=\lim_{n\to\infty}\int_0^1\left(a\cos2\pi x+b\sin2\pi x\right)dx \\
  &=0
\end{align}
より，
\begin{align}
  s_n&=\dfrac{1}{n}\sum_{k=0}^{n-1}L_k \\
  &=\dfrac{1}{n}\left\{n(1+a^2+b^2)-2\sum_{k=0}^{n-1}a_k\right\} \\
  &\lim_{n\to\infty}1+a^2+b^2
\end{align}
となって
\begin{align}
  \lambda_p=\lim{n\to\infty}s_n=1+a^2+b^2\cdots\text{(答)}
\end{align}
で，これは$P(0,0)$の時に最小であること明らかである．$\cdots$(答)

{\bf[解注]}$\sum{k=0}{n-1}a_k=0$は極限を取る前から成立していること明白だが，減点される恐れありと判断してあえて区分求積を用いた．
\newpage
\end{document}
